\section{Pure state quantum information}\label{sec:qubit_quantum_info}

Mathematically, a system of $n$ qubits is described by a $2^n$-dimensional complex vector space $\mathbb{C}^{2^n}$, which we refer to as the \emph{Hilbert space} $\HHH$ of the system.
Pure quantum states are equivalence classes of vectors in this space, under the equivalence $\ket{\psi} \sim \ket{\phi} \iff \ket{\phi} = \lambda \ket{\psi}$ for some non-zero $\lambda \in \mathbb{C}$.
This $\lambda$ is often called a \textit{global phase}.
A representative $\ket{\psi}$ is normalised if it has unit length.

\subsection{Evolution}

A pure quantum state evolves under the action of unitary operators and measurements.
A unitary operator is a linear map $U: \mathcal{H} \to \mathcal{H}$ satisfying $UU^\dagger = U^\dagger U = I$,
where $U^\dagger$ denotes the conjugate transpose.
If a system is in state $\ket{\psi}$ and we apply a unitary $U \in \mathcal{U}(2^n)$, the new state is $U\ket{\psi}$ -- no surprises here.

A \emph{measurement} on a quantum system is specified by a collection of linear operators $\{M_m\}_{m \in \mathcal{M}}$, where $\mathcal{M}$ is the set of possible measurement outcomes, satisfying:
\begin{equation}
    \sum_{m \in \mathcal{M}} M_m^\dagger M_m = I.
\end{equation}
Note that $\mathcal{M}$ can be any set -- the measurement outcome labels need not be real numbers (or even numbers at all!).
When we perform the measurement on a representative $\ket{\psi}$ of a state, the \textit{Born rule} says we obtain outcome $m \in \mathcal{M}$ with probability:
\begin{equation}\label{eq:born_rule}
    p(m | \psi) = \frac{\bra{\psi} M_m^\dagger M_m \ket{\psi}}{\braket{\psi | \psi}}
\end{equation}
If the representative is normalised, this is just $\bra{\psi} M_m^\dagger M_m \ket{\psi}$.
The post-measurement state after obtaining outcome $m$ is just $M_m \ket{\psi}$, i.e.\ the original state but with the $m$th operator $M_m$ applied.
For a normalised representative, the renormalised post-measurement state is:
\begin{equation}
    \ket{\psi_m} = \frac{M_m \ket{\psi}}{\sqrt{p(m | \psi)}}.
\end{equation}
A special case is a \emph{projective measurement}, where each $M_m$ is a projector: $M_m = \Pi_m$ with $\Pi_m^2 = \Pi_m = \Pi_m^\dagger$.
The completeness relation becomes $\sum_{m} \Pi_m = I$, and the projectors are mutually orthogonal: $\Pi_m \Pi_{m'} = \delta_{m,m'} \Pi_m$.

\subsection{The Pauli group}\label{subsec:pauli_group}

The single-qubit \emph{Pauli matrices} are
\begin{equation}
    I = \begin{pmatrix} 1 & 0 \\ 0 & 1 \end{pmatrix}, \quad
    X = \begin{pmatrix} 0 & 1 \\ 1 & 0 \end{pmatrix}, \quad
    Y = \begin{pmatrix} 0 & -i \\ i & 0 \end{pmatrix}, \quad
    Z = \begin{pmatrix} 1 & 0 \\ 0 & -1 \end{pmatrix}.
\end{equation}
These satisfy $X^2 = Y^2 = Z^2 = I$, $XYZ = iI$, and $PQ = -QP$ for distinct $Q, P \in \{X, Y, Z\}\}$.
The $n$-qubit \emph{Pauli group} $\mathcal{P}_n$ is the group generated by tensor products of single-qubit Paulis with phases:
\begin{equation}
    \mathcal{P}_n = \{ \pm 1, \pm i \} \times \{ I, X, Y, Z \}^{\otimes n}.
\end{equation}
More explicitly, elements of $\mathcal{P}_n$ are of the form $\omega (P_1 \otimes \cdots \otimes P_n)$ where $\omega \in \{1, i, -1, -i\}$ and each $P_j \in \{I, X, Y, Z\}$.
We will say a Pauli in the form above is in \textit{standard form} (it need not be -- e.g.\ we can write the Pauli above as $1 (\omega P_1 \otimes \ldots \otimes P_n)$, where $\omega P_1$ need not be in $\{I, X, Y, Z\}$).
By a small abuse of notation we will often use $I$ for both the single-qubit identity operator and the $n$-qubit identity operator.
More notation: given binary vectors $\boldsymbol{a}, \boldsymbol{b} \in \mathbb{F}_2^n$, we define $i^c X^{\bsym{a}}Z^{\bsym{b}}$ to be the Pauli $i^c X_1^{a_1} Z_1^{b_1} \ldots X_n^{a_n} Z_n^{a_n}$.
Given another two vectors $\boldsymbol{a}', \boldsymbol{b}' \in \mathbb{F}_2^n$, we define the \textit{symplectic product} of $\binom{\bsym{a}}{\bsym{b}}$ and $\binom{\bsym{a}'}{\bsym{b}'}$ to be:
\begin{equation}
    \tbinom{\bsym{a}}{\bsym{b}} \odot \tbinom{\bsym{a}'}{\bsym{b}'} = \bsym{a} \cdot \bsym{b}' - \bsym{b} \cdot \bsym{a}'
\end{equation}
where $\cdot$ denotes the standard dot product.

\subsection{Measuring a Pauli}

When we speak of \emph{measuring a Pauli operator} $P \in \mathcal{P}_n$, we mean performing a projective measurement in the eigenbasis of $P$, with measurement outcomes given by the corresponding eigenvalues.
Every Pauli is \emph{normal} (satisfies $P^\dagger P = P P^\dagger$), and normal operators are diagonalizable by the spectral theorem, so every Pauli can be measured.
Some sources in the literature ask for an operator to be Hermitian if we're going to measure it; this just ensures that the measurement outcomes (eigenvalues) are real.
Exactly half of the Paulis in $\PPP_n$ are Hermitian; these are those of the form $P = \omega (P_1 \otimes \ldots \otimes P_n)$ where $\omega = \pm 1$ and $P_j \in \{I, X, Y, Z\}$.
The remainder (where $\omega = \pm i$) are only normal.
But all of them can be measured!
Measuring non-Hermitian $P = \pm i (P_1 \otimes \ldots \otimes P_n)$ is equivalent to measuring Hermitian $P = \pm 1 (P_1 \otimes \ldots \otimes P_n)$; the only difference is that the outcomes are labelled differently, with $i$ instead of $1$ and $-i$ instead of $-1$.

Let $\Pi_{\lambda P}$ denote the projector $\frac{1}{2}(I + \lambda P)$ onto the $+1$-eigenspace of $\lambda P$; this is exactly the same thing as saying the projector onto the $\lambda$-eigenspace of $P$.
So the projector $\frac{1}{2}(I - \lambda P)$ onto the $-\lambda$-eigenspace can be written $\Pi_{-\lambda P}$.
These satisfy
\begin{equation}
    \Pi_{\lambda P} + \Pi_{-\lambda P} = I, \quad \Pi_{\lambda P} \Pi_{-\lambda P} = 0, \quad \Pi_{\pm \lambda P}^2 = \Pi_{\pm \lambda P}.
\end{equation}
The measurement is then defined as the operators $\{\Pi_{\lambda P}, \Pi_{-\lambda P}\}$, with outcomes labeled by the eigenvalues $\mathcal{M} = \{\lambda, -\lambda\}$ (or equivalently $\{1, -1\}$).

When we measure $P$ on a normalised state $\ket{\psi}$, we obtain outcome $m \in \{\lambda, -\lambda\}$ with probability $p(m | \psi) = \bra{\psi} \Pi_{mP} \ket{\psi}$, and the post-measurement state is
\begin{equation}
    \ket{\psi_m} = \frac{\Pi_{mP} \ket{\psi}}{\sqrt{p(m | \psi)}}.
\end{equation}
Note that $\ket{\psi_m}$ is a $+m$ eigenstate of $P$: we have $P \ket{\psi_m} = m \ket{\psi_m}$.

\subsection{The Clifford group}

The $n$-qubit \emph{Clifford group} $\mathcal{C}_n$ is defined as the normalizer of the Pauli group within the unitary group:
\begin{equation}
    \mathcal{C}_n = \{ U \in \mathcal{U}(2^n) : U \mathcal{P}_n U^\dagger = \mathcal{P}_n \}.
\end{equation}
In other words, Clifford unitaries are precisely those unitaries that map Pauli operators to Pauli operators under conjugation.

Common examples of single-qubit Clifford gates include the Hadamard gate
\begin{equation}
    H = \frac{1}{\sqrt{2}} \begin{pmatrix} 1 & 1 \\ 1 & -1 \end{pmatrix},
\end{equation}
the phase gate
\begin{equation}
    S = \begin{pmatrix} 1 & 0 \\ 0 & i \end{pmatrix}.
\end{equation}
and the controlled-NOT or CNOT gate
\begin{equation}
    \text{CNOT} = \begin{pmatrix} 1 & 0 & 0 & 0 \\ 0 & 1 & 0 & 0 \\ 0 & 0 & 0 & 1 \\ 0 & 0 & 1 & 0 \end{pmatrix}.
\end{equation}
In fact, the gates $\{H, S, \text{CNOT}\}$ generate the full Clifford group on any number of qubits.
